\begin{abstract}
Population prompts are widely used to generate synthetic survey responses, but they combine information supplied at inference with associations already encoded during pretraining. We introduce an alternative construction that maps declared aggregate preference anchors into group-indexed choice policies. For each population, the signs of six Global Preferences Survey (GPS) coordinates deterministically label a shared bank of paired synthetic responses, and Direct Preference Optimization fits a parameter-efficient adapter to those comparisons. We evaluate the adapters on candidate World Values Survey (WVS) items using prompts that omit country names and distinguish four questions: recovery of the imposed labels, transfer of the anchor signal to new text, coherence between the GPS anchors and human WVS responses, and agreement between adapter and human scores. The adapters recover the imposed pairwise labels. On a purposively selected sixteen-country development panel, adapter trust scores completely separate the two GPS-sign groups and have a rank correlation of (0.74) with continuous GPS trust scores. Human--GPS and adapter--human associations remain unresolved on the same panel, and results for the other preference dimensions are heterogeneous. These findings show that an anchored policy can retain a declared aggregate signal without thereby reproducing human response patterns. The contribution is therefore both an inspectable construction and an evaluation framework that separates anchor transfer from human criterion agreement.
\end{abstract}

\noindent\textbf{Keywords:} synthetic cultural agents; aggregate anchors; direct preference optimization; survey methodology; measurement validity.
